% GENERATED FILE. Edit canonical lessons, then run npm run generate:books.
% canonical-source-hash: fnv1a64:4dbf7668bd3397e1
% canonical-lessons: PA-C322-satt, PA-C322-zakham, PA-C322-sujjna, PA-C322-patti, PA-C322-malham, PA-R322-review-316-322

\chapter{A Cut and a Bruise}
\label{ch:pa-a2-a-cut-and-a-bruise}
\hlchaptermodality{322}

\begin{chapteropening}
\textbf{By the end of this chapter:} \emph{I can say where I got hurt and what I put on it.}

Complete the last lesson of chapter 322: I can say where I got hurt and what I put on it.
\end{chapteropening}

\section[\texorpdfstring{\pa{ਸੱਟ} (saṭṭ)}{ਸੱਟ (saṭṭ)}]{\pa{ਸੱਟ} (saṭṭ) — an injury}

\label{lesson:PA-C322-satt}



\begin{quote}
Before the new one: say the Punjabi for breath, then the Punjabi for to swallow.
\end{quote}

\subsection*{You'll want to know: \pa{ਸੱਟ}}
\textbf{\pa{ਸੱਟ}} — \emph{saṭṭ} — \textquotedblleft{}an injury\textquotedblright{}.

\begin{quote}
\pa{ਮੇਰੀ} \pa{ਲੱਤ} \pa{ਉੱਤੇ} \pa{ਇੱਕ} \pa{ਵੱਡੀ} \pa{ਸੱਟ} \pa{ਹੈ।}
\end{quote}

\begin{quote}
\emph{merī latt utte ikk vaḍḍī saṭṭ hai.}
\end{quote}

\begin{quote}
I have a big injury on my leg.
\end{quote}

A feminine noun.

\begin{cousinweb}
\textbf{\pa{ਸੱਟ}} (\emph{saṭṭ}) is inherited from Sanskrit \emph{sṛṣṭi}, \textquotedblleft{}emission, creation\textquotedblright{}; CDIAL lists Punjabi \emph{saṭṭ}, \textquotedblleft{}a blow, a wound\textquotedblright{}, under it (CDIAL 13579). How the sense moved is not explained there.
\end{cousinweb}

\begin{grammarlens}[title={What you've built}]
One more word for this chapter.
\end{grammarlens}

\subsection*{Your turn}
\begin{itemize}
\raggedright
  \item \emph{Say it:} \emph{saṭṭ}
  \item \emph{Say it:} \emph{saṭṭ}, once more
  \item \emph{Say it:} \emph{merī latt utte ikk vaḍḍī saṭṭ hai.} — I have a big injury on my leg.
  \item \emph{Recall:} say the Punjabi for breath, then the Punjabi for to swallow, then say \emph{saṭṭ} again
\end{itemize}

\begin{tcolorbox}[breakable,colback=teal!4,colframe=teal!35!black,title={Before you move on}]
What does \pa{ਸੱਟ} mean? (\textquotedblleft{}An injury\textquotedblright{}.) Say it once more.
\end{tcolorbox}
\section[\texorpdfstring{\pa{ਜ਼ਖ਼ਮ} (zakham)}{ਜ਼ਖ਼ਮ (zakham)}]{\pa{ਜ਼ਖ਼ਮ} (zakham) — a wound}

\label{lesson:PA-C322-zakham}



\begin{quote}
Before the new one: say the Punjabi for to swallow, then the Punjabi for an injury.
\end{quote}

\subsection*{You'll want to know: \pa{ਜ਼ਖ਼ਮ}}
\textbf{\pa{ਜ਼ਖ਼ਮ}} — \emph{zakham} — \textquotedblleft{}a wound\textquotedblright{}.

\begin{quote}
\pa{ਚਾਕੂ} \pa{ਨਾਲ} \pa{ਮੇਰੀ} \pa{ਉਂਗਲ} \pa{ਉੱਤੇ} \pa{ਜ਼ਖ਼ਮ} \pa{ਹੋਇਆ।}
\end{quote}

\begin{quote}
\emph{chākū nāl merī uṅgal utte zakham hoiā.}
\end{quote}

\begin{quote}
I got a cut on my finger from the knife.
\end{quote}

A masculine noun.

\begin{cousinweb}
\textbf{\pa{ਜ਼ਖ਼ਮ}} (\emph{zakham}) is a loan, usually traced to Persian \emph{zakhm}, \textquotedblleft{}a wound\textquotedblright{}.
\end{cousinweb}

\begin{grammarlens}[title={What you've built}]
One more word for this chapter.
\end{grammarlens}

\subsection*{Your turn}
\begin{itemize}
\raggedright
  \item \emph{Say it:} \emph{zakham}
  \item \emph{Say it:} \emph{zakham}, once more
  \item \emph{Say it:} \emph{chākū nāl merī uṅgal utte zakham hoiā.} — I got a cut on my finger from the knife.
  \item \emph{Recall:} say the Punjabi for to swallow, then the Punjabi for an injury, then say \emph{zakham} again
  \item \emph{Write it:} copy \textbf{\pa{ਜ਼ਖ਼ਮ}} once, checking the whole word
\end{itemize}

\begin{tcolorbox}[breakable,colback=teal!4,colframe=teal!35!black,title={Before you move on}]
What does \pa{ਜ਼ਖ਼ਮ} mean? (\textquotedblleft{}A wound\textquotedblright{}.) Say it once more.
\end{tcolorbox}
\section[\texorpdfstring{\pa{ਸੁੱਜਣਾ} (sujjṇā)}{ਸੁੱਜਣਾ (sujjṇā)}]{\pa{ਸੁੱਜਣਾ} (sujjṇā) — to swell up}

\label{lesson:PA-C322-sujjna}



\begin{quote}
Before the new one: say the Punjabi for an injury, then the Punjabi for a wound.
\end{quote}

\subsection*{You'll want to know: \pa{ਸੁੱਜਣਾ}}
\textbf{\pa{ਸੁੱਜਣਾ}} — \emph{sujjṇā} — \textquotedblleft{}to swell up\textquotedblright{}.

\begin{quote}
\pa{ਸੱਟ} \pa{ਤੋਂ} \pa{ਬਾਅਦ} \pa{ਮੇਰਾ} \pa{ਗਿੱਟਾ} \pa{ਸੁੱਜਿਆ।}
\end{quote}

\begin{quote}
\emph{saṭṭ ton bāad merā giṭṭā sujjiā.}
\end{quote}

\begin{quote}
After the injury my ankle swelled up.
\end{quote}

Past (it …): \pa{ਉਹ} \pa{ਗੋਡਾ} \pa{ਬਹੁਤ} \pa{ਸੁੱਜਿਆ}, \emph{oh goḍā bahut sujjiā}.

\begin{cousinweb}
\textbf{\pa{ਸੁੱਜਣਾ}} (\emph{sujjṇā}) is inherited from Sanskrit \emph{śūyate}, \textquotedblleft{}is swollen\textquotedblright{}; CDIAL lists Punjabi \emph{sujjṇā}, \textquotedblleft{}to swell\textquotedblright{}, under it (CDIAL 12568).
\end{cousinweb}

\begin{grammarlens}[title={What you've built}]
One more word for this chapter.
\end{grammarlens}

\subsection*{Your turn}
\begin{itemize}
\raggedright
  \item \emph{Say it:} \emph{sujjṇā}
  \item \emph{Say it:} \emph{sujjṇā}, once more
  \item \emph{Say it:} \emph{saṭṭ ton bāad merā giṭṭā sujjiā.} — After the injury my ankle swelled up.
  \item \emph{Recall:} say the Punjabi for an injury, then the Punjabi for a wound, then say \emph{sujjṇā} again
  \item \emph{Write it:} copy \textbf{\pa{ਸੁੱਜਣਾ}} once, checking the whole word
\end{itemize}

\begin{tcolorbox}[breakable,colback=teal!4,colframe=teal!35!black,title={Before you move on}]
What does \pa{ਸੁੱਜਣਾ} mean? (\textquotedblleft{}To swell up\textquotedblright{}.) Say it once more.
\end{tcolorbox}
\section[\texorpdfstring{\pa{ਪੱਟੀ} (paṭṭī)}{ਪੱਟੀ (paṭṭī)}]{\pa{ਪੱਟੀ} (paṭṭī) — a bandage}

\label{lesson:PA-C322-patti}



\begin{quote}
Before the new one: say the Punjabi for a wound, then the Punjabi for to swell up.
\end{quote}

\subsection*{You'll want to know: \pa{ਪੱਟੀ}}
\textbf{\pa{ਪੱਟੀ}} — \emph{paṭṭī} — \textquotedblleft{}a bandage\textquotedblright{}.

\begin{quote}
\pa{ਮੈਂ} \pa{ਜ਼ਖ਼ਮ} \pa{ਉੱਤੇ} \pa{ਸਾਫ਼} \pa{ਪੱਟੀ} \pa{ਬੰਨ੍ਹੀ।}
\end{quote}

\begin{quote}
\emph{maiṁ zakham utte sāf paṭṭī bannhī.}
\end{quote}

\begin{quote}
I tied a clean bandage over the wound.
\end{quote}

A feminine noun.

\begin{cousinweb}
\textbf{\pa{ਪੱਟੀ}} (\emph{paṭṭī}) is inherited from Sanskrit \emph{paṭṭa}, \textquotedblleft{}cloth; a bandage\textquotedblright{}; CDIAL lists Punjabi \emph{paṭṭī}, \textquotedblleft{}coarse woollen cloth, bandage\textquotedblright{}, under it (CDIAL 7700).
\end{cousinweb}

\begin{grammarlens}[title={What you've built}]
One more word for this chapter.
\end{grammarlens}

\subsection*{Your turn}
\begin{itemize}
\raggedright
  \item \emph{Say it:} \emph{paṭṭī}
  \item \emph{Say it:} \emph{paṭṭī}, once more
  \item \emph{Say it:} \emph{maiṁ zakham utte sāf paṭṭī bannhī.} — I tied a clean bandage over the wound.
  \item \emph{Recall:} say the Punjabi for a wound, then the Punjabi for to swell up, then say \emph{paṭṭī} again
  \item \emph{Write it:} copy \textbf{\pa{ਪੱਟੀ}} once, checking the whole word
\end{itemize}

\begin{tcolorbox}[breakable,colback=teal!4,colframe=teal!35!black,title={Before you move on}]
What does \pa{ਪੱਟੀ} mean? (\textquotedblleft{}A bandage\textquotedblright{}.) Say it once more.
\end{tcolorbox}
\section[\texorpdfstring{\pa{ਮਲ੍ਹਮ} (malham)}{ਮਲ੍ਹਮ (malham)}]{\pa{ਮਲ੍ਹਮ} (malham) — an ointment}

\label{lesson:PA-C322-malham}



\begin{quote}
Before the new one: say the Punjabi for to swell up, then the Punjabi for a bandage.
\end{quote}

\subsection*{You'll want to know: \pa{ਮਲ੍ਹਮ}}
\textbf{\pa{ਮਲ੍ਹਮ}} — \emph{malham} — \textquotedblleft{}an ointment\textquotedblright{}.

\begin{quote}
\pa{ਰਾਤ} \pa{ਨੂੰ} \pa{ਮੈਂ} \pa{ਸੱਟ} \pa{ਉੱਤੇ} \pa{ਮਲ੍ਹਮ} \pa{ਮਲੀ।}
\end{quote}

\begin{quote}
\emph{rāt nū̃ maiṁ saṭṭ utte malham malī.}
\end{quote}

\begin{quote}
At night I rubbed ointment on the injury.
\end{quote}

A feminine noun.

\begin{cousinweb}
\textbf{\pa{ਮਲ੍ਹਮ}} (\emph{malham}) is a loan, usually traced through Persian to Arabic \emph{marham}, \textquotedblleft{}an ointment\textquotedblright{}.
\end{cousinweb}

\begin{grammarlens}[title={What you've built}]
One more word for this chapter.
\end{grammarlens}

\subsection*{Your turn}
\begin{itemize}
\raggedright
  \item \emph{Say it:} \emph{malham}
  \item \emph{Say it:} \emph{malham}, once more
  \item \emph{Say it:} \emph{rāt nū̃ maiṁ saṭṭ utte malham malī.} — At night I rubbed ointment on the injury.
  \item \emph{Recall:} say the Punjabi for to swell up, then the Punjabi for a bandage, then say \emph{malham} again
  \item \emph{Write it:} copy \textbf{\pa{ਮਲ੍ਹਮ}} once, checking the whole word
\end{itemize}

\begin{tcolorbox}[breakable,colback=teal!4,colframe=teal!35!black,title={Before you move on}]
What does \pa{ਮਲ੍ਹਮ} mean? (\textquotedblleft{}An ointment\textquotedblright{}.) Say it once more.
\end{tcolorbox}
\section[(dialogue)]{Looking back: from Trouble in the Lane to A Cut and a Bruise}

\label{lesson:PA-R322-review-316-322}



\begin{quote}
Say the Punjabi for to fight and a quarrel.
\end{quote}

\subsection*{Your turn}
Say these aloud:

\begin{itemize}
\raggedright
  \item to fight, a quarrel, to hit, a swear word, a reconciliation
  \item a particular time, to go by (of time), an era, to regret, to spend time
  \item to fly, to burst, a collision, to sink, a boatman
  \item how someone is feeling, how things stand, health, healthy, well again after an illness
  \item a doctor, a patient, medical treatment, the pulse, an injection
  \item a head cold, a cough, a sneeze, breath, to swallow
  \item an injury, a wound, to swell up, a bandage, an ointment
\end{itemize}

\begin{tcolorbox}[breakable,colback=teal!4,colframe=teal!35!black,title={Before you move on}]
To fight? (\textbf{\pa{ਲੜਨਾ}}, \emph{laṛnā}.) A quarrel? (\textbf{\pa{ਝਗੜਾ}}, \emph{jhagṛā}.) To hit? (\textbf{\pa{ਮਾਰਨਾ}}, \emph{mārnā}.) A swear word? (\textbf{\pa{ਗਾਲ੍ਹ}}, \emph{gālh}.) A reconciliation? (\textbf{\pa{ਸੁਲ੍ਹਾ}}, \emph{sulhā}.) A particular time? (\textbf{\pa{ਵੇਲਾ}}, \emph{velā}.) To go by (of time)? (\textbf{\pa{ਬੀਤਣਾ}}, \emph{bītṇā}.) An era? (\textbf{\pa{ਜ਼ਮਾਨਾ}}, \emph{zamānā}.) To regret? (\textbf{\pa{ਪਛਤਾਉਣਾ}}, \emph{pachhtāuṇā}.) To spend time? (\textbf{\pa{ਗੁਜ਼ਾਰਨਾ}}, \emph{guzārnā}.) To fly? (\textbf{\pa{ਉੱਡਣਾ}}, \emph{uḍḍṇā}.) To burst? (\textbf{\pa{ਫੁੱਟਣਾ}}, \emph{phuṭṭṇā}.) A collision? (\textbf{\pa{ਟੱਕਰ}}, \emph{ṭakkar}.) To sink? (\textbf{\pa{ਡੁੱਬਣਾ}}, \emph{ḍubbṇā}.) A boatman? (\textbf{\pa{ਮਲਾਹ}}, \emph{malāh}.) How someone is feeling? (\textbf{\pa{ਤਬੀਅਤ}}, \emph{tabīat}.) How things stand? (\textbf{\pa{ਹਾਲ}}, \emph{hāl}.) Health? (\textbf{\pa{ਸਿਹਤ}}, \emph{sihat}.) Healthy? (\textbf{\pa{ਤੰਦਰੁਸਤ}}, \emph{tandrusat}.) Well again after an illness? (\textbf{\pa{ਰਾਜ਼ੀ}}, \emph{rāzī}.) A doctor? (\textbf{\pa{ਡਾਕਟਰ}}, \emph{ḍākṭar}.) A patient? (\textbf{\pa{ਮਰੀਜ਼}}, \emph{marīz}.) Medical treatment? (\textbf{\pa{ਇਲਾਜ}}, \emph{ilāj}.) The pulse? (\textbf{\pa{ਨਬਜ਼}}, \emph{nabaz}.) An injection? (\textbf{\pa{ਟੀਕਾ}}, \emph{ṭīkā}.) A head cold? (\textbf{\pa{ਜ਼ੁਕਾਮ}}, \emph{zukām}.) A cough? (\textbf{\pa{ਖੰਘ}}, \emph{khaṅgh}.) A sneeze? (\textbf{\pa{ਛਿੱਕ}}, \emph{chhikk}.) Breath? (\textbf{\pa{ਸਾਹ}}, \emph{sāh}.) To swallow? (\textbf{\pa{ਨਿਗਲਣਾ}}, \emph{nigalṇā}.) An injury? (\textbf{\pa{ਸੱਟ}}, \emph{saṭṭ}.) A wound? (\textbf{\pa{ਜ਼ਖ਼ਮ}}, \emph{zakham}.) To swell up? (\textbf{\pa{ਸੁੱਜਣਾ}}, \emph{sujjṇā}.) A bandage? (\textbf{\pa{ਪੱਟੀ}}, \emph{paṭṭī}.) An ointment? (\textbf{\pa{ਮਲ੍ਹਮ}}, \emph{malham}.)
\end{tcolorbox}
